% GENERATED FILE. Edit canonical lessons, then run npm run generate:books.
% canonical-source-hash: fnv1a64:39fbe68165e98ff5
% canonical-lessons: SA-C298-ksatih, SA-C298-nasayati, SA-C298-bhagnah, SA-C298-paurah, SA-C298-sesah

\chapter{After the Fire}
\label{ch:sa-after-the-fire}
\hlchaptermodality{298}

\begin{chapteropening}
\textbf{By the end of this chapter:} \emph{I can describe the damage a fire left behind.}

Complete the last lesson of chapter 298: I can describe the damage a fire left behind.
\end{chapteropening}

\section[\texorpdfstring{\sk{क्षतिः} (kṣatiḥ)}{क्षतिः (kṣatiḥ)}]{\sk{क्षतिः} (kṣatiḥ) — damage; harm}

\label{lesson:SA-C298-ksatih}



\begin{quote}
Before the new one: say the Sanskrit for puts out a fire, then the Sanskrit for dies down.
\end{quote}

\subsection*{You'll want to know: \sk{क्षतिः}}
\textbf{\sk{क्षतिः}} — \emph{kṣatiḥ} — \textquotedblleft{}damage; harm\textquotedblright{}.

\begin{quote}
\sk{अग्निना} \sk{ग्रामस्य} \sk{क्षतिः} \sk{भवति।}
\end{quote}

\begin{quote}
\emph{agninā grāmasya kṣatiḥ bhavati.}
\end{quote}

\begin{quote}
The fire does damage to the village.
\end{quote}

A feminine noun.

\begin{cousinweb}
\textbf{\sk{क्षतिः}} (\emph{kṣatiḥ}) MW gives \emph{kṣati} as injury, loss and damage, beside \emph{kṣata}, \textquotedblleft{}wounded\textquotedblright{}, but gives no root in the entry, so its origin is left uncertain here (MW).
\end{cousinweb}

\begin{grammarlens}[title={What you've built}]
One more word for telling what happened.
\end{grammarlens}

\subsection*{Your turn}
\begin{itemize}
\raggedright
  \item \emph{Say it:} \emph{kṣatiḥ}
  \item \emph{Say it:} \emph{kṣatiḥ}, once more
  \item \emph{Say it:} \emph{agninā grāmasya kṣatiḥ bhavati.} — The fire does damage to the village.
  \item \emph{Recall:} say the Sanskrit for puts out a fire, then the Sanskrit for dies down, then say \emph{kṣatiḥ} again
\end{itemize}

\begin{tcolorbox}[breakable,colback=teal!4,colframe=teal!35!black,title={Before you move on}]
What does \sk{क्षतिः} mean? (\textquotedblleft{}Damage; harm\textquotedblright{}.) Say it once more.
\end{tcolorbox}
\section[\texorpdfstring{\sk{नाशयति} (nāśayati)}{नाशयति (nāśayati)}]{\sk{नाशयति} (nāśayati) — destroys}

\label{lesson:SA-C298-nasayati}



\begin{quote}
Before the new one: say the Sanskrit for dies down, then the Sanskrit for damage.
\end{quote}

\subsection*{You'll want to know: \sk{नाशयति}}
\textbf{\sk{नाशयति}} — \emph{nāśayati} — \textquotedblleft{}destroys\textquotedblright{}.

\begin{quote}
\sk{अग्निः} \sk{गृहाणि} \sk{नाशयति।}
\end{quote}

\begin{quote}
\emph{agniḥ gṛhāṇi nāśayati.}
\end{quote}

\begin{quote}
The fire destroys the houses.
\end{quote}

Past (it …): \sk{अग्निः} \sk{गृहाणि} \sk{नाशितवान्}, \emph{agniḥ gṛhāṇi nāśitavān}.

\begin{cousinweb}
\textbf{\sk{नाशयति}} (\emph{nāśayati}) comes from the root √naś, \textquotedblleft{}to be lost, perish\textquotedblright{}, which also gives \emph{naśyati}, in its causative form: to make perish, destroy (MW).
\end{cousinweb}

\begin{grammarlens}[title={What you've built}]
One more word for telling what happened.
\end{grammarlens}

\subsection*{Your turn}
\begin{itemize}
\raggedright
  \item \emph{Say it:} \emph{nāśayati}
  \item \emph{Say it:} \emph{nāśayati}, once more
  \item \emph{Say it:} \emph{agniḥ gṛhāṇi nāśayati.} — The fire destroys the houses.
  \item \emph{Recall:} say the Sanskrit for dies down, then the Sanskrit for damage, then say \emph{nāśayati} again
\end{itemize}

\begin{tcolorbox}[breakable,colback=teal!4,colframe=teal!35!black,title={Before you move on}]
What does \sk{नाशयति} mean? (\textquotedblleft{}Destroys\textquotedblright{}.) Say it once more.
\end{tcolorbox}
\section[\texorpdfstring{\sk{भग्नः} (bhagnaḥ)}{भग्नः (bhagnaḥ)}]{\sk{भग्नः} (bhagnaḥ) — broken}

\label{lesson:SA-C298-bhagnah}



\begin{quote}
Before the new one: say the Sanskrit for damage, then the Sanskrit for destroys.
\end{quote}

\subsection*{You'll want to know: \sk{भग्नः}}
\textbf{\sk{भग्नः}} — \emph{bhagnaḥ} — \textquotedblleft{}broken\textquotedblright{}.

\begin{quote}
\sk{गृहस्य} \sk{छदिः} \sk{भग्नम्} \sk{अस्ति।}
\end{quote}

\begin{quote}
\emph{gṛhasya chadiḥ bhagnam asti.}
\end{quote}

\begin{quote}
The roof of the house is broken.
\end{quote}

\begin{cousinweb}
\textbf{\sk{भग्नः}} (\emph{bhagnaḥ}) comes from the root √bhañj, \textquotedblleft{}to break, shatter\textquotedblright{}; MW gives \emph{bhagna} as broken, shattered, split (MW).
\end{cousinweb}

\begin{grammarlens}[title={What you've built}]
One more word for telling what happened.
\end{grammarlens}

\subsection*{Your turn}
\begin{itemize}
\raggedright
  \item \emph{Say it:} \emph{bhagnaḥ}
  \item \emph{Say it:} \emph{bhagnaḥ}, once more
  \item \emph{Say it:} \emph{gṛhasya chadiḥ bhagnam asti.} — The roof of the house is broken.
  \item \emph{Recall:} say the Sanskrit for damage, then the Sanskrit for destroys, then say \emph{bhagnaḥ} again
\end{itemize}

\begin{tcolorbox}[breakable,colback=teal!4,colframe=teal!35!black,title={Before you move on}]
What does \sk{भग्नः} mean? (\textquotedblleft{}Broken\textquotedblright{}.) Say it once more.
\end{tcolorbox}
\section[\texorpdfstring{\sk{पौरः} (pauraḥ)}{पौरः (pauraḥ)}]{\sk{पौरः} (pauraḥ) — a citizen; a townsman}

\label{lesson:SA-C298-paurah}



\begin{quote}
Before the new one: say the Sanskrit for destroys, then the Sanskrit for broken.
\end{quote}

\subsection*{You'll want to know: \sk{पौरः}}
\textbf{\sk{पौरः}} — \emph{pauraḥ} — \textquotedblleft{}a citizen; a townsman\textquotedblright{}.

\begin{quote}
\sk{आहताः} \sk{पौराः} \sk{औषधं} \sk{पिबन्ति।}
\end{quote}

\begin{quote}
\emph{āhatāḥ paurāḥ auṣadhaṃ pibanti.}
\end{quote}

\begin{quote}
The injured townspeople take medicine.
\end{quote}

A masculine noun.

\begin{cousinweb}
\textbf{\sk{पौरः}} (\emph{pauraḥ}) is built on \emph{pura}, \textquotedblleft{}a town\textquotedblright{}, with the lengthened first vowel that makes \textquotedblleft{}belonging to\textquotedblright{}; MW gives \emph{paura} as a townsman, a citizen (MW).
\end{cousinweb}

\begin{grammarlens}[title={What you've built}]
One more word for telling what happened.
\end{grammarlens}

\subsection*{Your turn}
\begin{itemize}
\raggedright
  \item \emph{Say it:} \emph{pauraḥ}
  \item \emph{Say it:} \emph{pauraḥ}, once more
  \item \emph{Say it:} \emph{āhatāḥ paurāḥ auṣadhaṃ pibanti.} — The injured townspeople take medicine.
  \item \emph{Recall:} say the Sanskrit for destroys, then the Sanskrit for broken, then say \emph{pauraḥ} again
\end{itemize}

\begin{tcolorbox}[breakable,colback=teal!4,colframe=teal!35!black,title={Before you move on}]
What does \sk{पौरः} mean? (\textquotedblleft{}A citizen; a townsman\textquotedblright{}.) Say it once more.
\end{tcolorbox}
\section[\texorpdfstring{\sk{शेषः} (śeṣaḥ)}{शेषः (śeṣaḥ)}]{\sk{शेषः} (śeṣaḥ) — what is left; the remainder}

\label{lesson:SA-C298-sesah}



\begin{quote}
Before the new one: say the Sanskrit for broken, then the Sanskrit for a citizen.
\end{quote}

\subsection*{You'll want to know: \sk{शेषः}}
\textbf{\sk{शेषः}} — \emph{śeṣaḥ} — \textquotedblleft{}what is left; the remainder\textquotedblright{}.

\begin{quote}
\sk{गृहस्य} \sk{शेषः} \sk{केवलं} \sk{भस्म} \sk{अस्ति।}
\end{quote}

\begin{quote}
\emph{gṛhasya śeṣaḥ kevalaṃ bhasma asti.}
\end{quote}

\begin{quote}
All that is left of the house is ash.
\end{quote}

A masculine noun.

\begin{cousinweb}
\textbf{\sk{शेषः}} (\emph{śeṣaḥ}) comes from the root √śiṣ, \textquotedblleft{}to leave, leave remaining\textquotedblright{}; MW gives \emph{śeṣa} as the remainder, what is left (MW).
\end{cousinweb}

\begin{grammarlens}[title={What you've built}]
One more word for telling what happened.
\end{grammarlens}

\subsection*{Your turn}
\begin{itemize}
\raggedright
  \item \emph{Say it:} \emph{śeṣaḥ}
  \item \emph{Say it:} \emph{śeṣaḥ}, once more
  \item \emph{Say it:} \emph{gṛhasya śeṣaḥ kevalaṃ bhasma asti.} — All that is left of the house is ash.
  \item \emph{Recall:} say the Sanskrit for broken, then the Sanskrit for a citizen, then say \emph{śeṣaḥ} again
\end{itemize}

\begin{tcolorbox}[breakable,colback=teal!4,colframe=teal!35!black,title={Before you move on}]
What does \sk{शेषः} mean? (\textquotedblleft{}What is left; the remainder\textquotedblright{}.) Say it once more.
\end{tcolorbox}
